\begin{definition}[Secant Line]
\label{def:secant-line}
Let $f : A \to \mathbb{R}$, and let $x_1,x_2 \in A$ with $x_1 \neq x_2$. The \emph{secant line} determined by the points
\[
(x_1,f(x_1))
\qquad \text{and} \qquad
(x_2,f(x_2))
\]
is the unique line passing through these two points. Its slope is
\[
\frac{f(x_2)-f(x_1)}{x_2-x_1}.
\]
\end{definition}

\begin{definition}[Tangent Line]
\label{def:tangent-line}
Let $f : A \to \mathbb{R}$, let $c \in A$, and suppose that $f$ is differentiable at $c$. The \emph{tangent line} to the graph of $f$ at the point $(c,f(c))$ is the line through $(c,f(c))$ with slope $f'(c)$. Its equation is
\[
y - f(c) = f'(c)(x-c).
\]
Equivalently,
\[
y = f(c) + f'(c)(x-c).
\]
\end{definition}




\begin{remark*}[Increment form]
With increments $\Delta x := x - c$ and $\Delta y := f(x) - f(c)$, the
difference quotient gives the exact relation
\[
  \Delta y = m_{\mathrm{sec}}(c,x) \cdot \Delta x
  = \frac{f(x)-f(c)}{x-c} \cdot (x-c).
\]
This identity is purely algebraic: it expresses the output change as
secant slope times input change. No limit is involved.
\end{remark*}

\begin{remark*}[Geometric interpretation]
The derivative $f'(c)$, when it exists, is the limit of the secant slopes
$m_{\mathrm{sec}}(c,x)$ as $x \to c$. Geometrically, the secant line through
$(c, f(c))$ and $(x, f(x))$ rotates toward the tangent line at $(c, f(c))$
as $x \to c$. The derivative is the slope of that limiting tangent.
See \hyperref[fig:secant-tangent]{Figure: Secant converging to tangent}.
\end{remark*}

% Figure: secant and tangent lines — extracted to figure file
\input{volume-iii/analysis/differentiation/notes/secant-tangent/figure1}